\chapter{La cascade dirigée}
\label{chap:cascade}

\begin{cle}
\textbf{Idée directrice.} La dépendance entre les cinq piliers n'est pas ajoutée après coup par
une copule : elle est \emph{produite} par la propagation d'un incident le long d'une matrice
d'influence dirigée. La cascade est ainsi la vue structurelle (bottom-up) de la
dépendance, en regard de la vue par copule (top-down). Le fil du
chapitre part du Vasicek à un facteur, l'ouvre en deux temps (systémique par pilier, puis
propagation dirigée), et remplace le seuil de défaut du crédit par un seuil ancré sur la barre
de matérialité du règlement.
\end{cle}

\section{Le point de départ : le Vasicek à un facteur et ses limites}

Le modèle structurel de Merton-Vasicek associe à chaque entité une variable latente. On retient
la convention où $X$ mesure le stress et non la santé (un $X$ élevé est une mauvaise
nouvelle), de sorte qu'une contagion positive encode bien une contamination :
\begin{equation}
X_i \;=\; \sqrt{\rho}\;Y \;+\; \sqrt{1-\rho}\;\varepsilon_i, \qquad \mathrm{Var}(X_i)=1,
\end{equation}
où $Y\sim\mathcal N(0,1)$ est le facteur systémique commun, $\varepsilon_i\sim\mathcal N(0,1)$ le
choc idiosyncratique, et $\rho\in[0,1]$ la part du systémique, commune à toutes les entités.
L'incident survient lorsque le stress franchit un seuil, $D_i=\mathbf 1\{X_i\ge K_i\}$, avec dans
le modèle classique $K_i=\Phi^{-1}(1-\mathrm{PD}_i)$.

\begin{attention}
\textbf{Deux hypothèses de ce modèle sont fausses pour le cyber sous DORA.} D'abord, la
dépendance y est symétrique (un seul $Y$, un seul $\rho$) : le modèle ne peut pas dire
\og P1 entraîne P2 plus que l'inverse\fg{}, alors que les données d'expert mobilisées ici
sont dirigées à $81\,\%$. Ensuite, le seuil $\Phi^{-1}(\mathrm{PD})$ suppose une PD stable, mesurable et
une valeur de firme sous-jacente, dont rien n'existe en cyber. Les deux généralisations qui
suivent lèvent ces deux limites.
\end{attention}

% SOURCES-SCRIPTS: 01 03

\section{Généralisation A : la structure systémique dépend du pilier}

On indexe désormais par entité~$i$ et pilier~$j\in\{1,\dots,5\}$, et l'on remplace le facteur
unique par un facteur systémique par pilier $Y_j$ et une sensibilité propre
$\rho_{ij}$ :
\begin{equation}
X_{ij} \;=\; \sqrt{\rho_{ij}}\;Y_j \;+\; \sqrt{1-\rho_{ij}}\;\varepsilon_{ij}.
\end{equation}
Chaque pilier a sa \og météo systémique\fg{} propre (l'état de la menace sur la chaîne
d'approvisionnement pour P4, l'état réglementaire pour P1), et les $Y_j$ sont corrélés entre eux
par une matrice $\Sigma_Y$ : c'est là que vivent proprement les dépendances croisées, au niveau
systémique. Sur des données cyber rares, on ne calibre pas un $\rho_{ij}$ par entité (environ
$5N$ paramètres, non identifiable) : on retient le compromis hiérarchique
$\rho_{ij}=\rho_j+u_{ij}$, où la sensibilité de chaque entité gravite autour de la moyenne de
son pilier $\rho_j$ (pooling partiel).

\begin{cle}
\textbf{La partie A n'est pas une contribution.} Le Merton-Vasicek multifacteur à
matrice de corrélation générale est celui de \citet{Li2000}, employé tel quel en risque de
crédit de portefeuille. La contribution propre ne commence qu'en partie~B, avec
la contagion dirigée. Le revendiquer clairement évite de se voir opposer une
littérature qu'on aurait ignorée.
\end{cle}

\section{Généralisation B : la propagation dirigée}

\begin{cle}
\textbf{Convention d'indices, à fixer une fois pour toutes.} Dans tout le mémoire,
\[
  \boxed{\;W_{jk}\;=\;\text{propension du pilier } j \text{ (source) à entraîner la
  défaillance du pilier } k \text{ (cible).}\;}
\]
La ligne $j$ décrit donc ce que $j$ \emph{émet}, et la colonne $k$ ce que $k$
\emph{reçoit}. C'est la convention du graphe d'arêtes vivantes (définition~\ref{def:cascade} :
l'arc $j\to k$ existe si $B_{jk}=1$) et celle de toutes les implémentations. Elle donne son sens
à la somme de ligne $s_j$, qui mesure l'émission du pilier $j$, et à la revendication
centrale \og{}P1 émet fort et reçoit peu\fg{} : P1 a la plus grande somme de ligne ($2{,}6$) et
une somme de colonne trois fois inférieure à celle de P2. Toute lecture inverse rendrait cette
revendication fausse.
\end{cle}

Pour encoder l'asymétrie, on ajoute un terme de réseau dirigé : le stress de $j$ subit celui des
autres piliers de la même entité, pondéré par une matrice asymétrique $W$. En respectant
la convention ci-dessus, la source est le premier indice, donc le stress \emph{reçu} par $j$
somme les colonnes :
\begin{equation}
X_{ij} \;=\; \underbrace{\sum_{k\neq j} W_{kj}\,X_{ik}}_{\text{contagion dirigée, } k \to j}
   \;+\; \sqrt{\rho_{ij}}\,Y_j \;+\; \sqrt{1-\rho_{ij}}\,\varepsilon_{ij},
\end{equation}
d'où, en forme matricielle sur les cinq piliers, la forme réduite
$\mathbf X_i=(I-W^{\top})^{-1}(\text{systémique}+\text{idio})$, \emph{sous réserve} que
$\rho(W)<1$, le rayon spectral étant invariant par transposition. Si
$W=0$ on retombe sur la partie~A, et avec un seul facteur sur le Vasicek standard :
l'emboîtement est propre.

\begin{figure}[htbp]\centering
\figover{H1_reseau_W.png}
\caption{Structure dirigée de $W$ sur les cinq piliers (P1 émet fort et reçoit peu) et
propagation d'un choc sur un sous-réseau.}
\end{figure}

% SOURCES-SCRIPTS: 29

\section{La normalisation de Leontief : une stabilité garantie}

\begin{attention}
\textbf{Brancher la matrice d'expert TRANS telle quelle est une erreur de catégorie.}
$(I-W)^{-1}$ ne représente une cascade que si la série $I+W+W^2+\cdots$ converge, soit
$\rho(W)<1$. Or $\rho(\mathrm{TRANS})=1{,}461>1$, et $(I-\mathrm{TRANS})^{-1}$ contient
21 entrées négatives sur 25 : un choc positif produirait un stress négatif ailleurs.
TRANS n'a jamais été une matrice de parts.
\end{attention}

La discipline vient du modèle entrées-sorties de Leontief, dont l'inverse
$B=(I-A)^{-1}=\sum_{k\ge0}A^k$ n'explose jamais parce que $A$ est une matrice de parts :
les colonnes somment à moins de un, la différence étant la valeur ajoutée, strictement positive.
On applique la même règle en lisant la ligne $j$ de $W$ comme les parts du stress que $j$
\emph{émet} vers les autres piliers :
\begin{equation}
\label{eq:norm}
\boxed{\;W \;=\; g \cdot \frac{\mathrm{TRANS}}{\max_j \sum_k \mathrm{TRANS}_{jk}}\;,
\qquad g \in (0,\,1]\;}
\end{equation}
où $g$ est la part de contagion du pilier le plus exposé. La contrainte $g\le1$ n'est pas ad
hoc : $W$ est positive et chaque ligne somme à au plus $g$, donc par Perron-Frobenius (le strict
analogue de la valeur ajoutée positive de Leontief)
\begin{equation}
\rho(W)\;\le\;\max_j\sum_k W_{jk}\;\le\;g\;<\;1,
\end{equation}
et $(I-W)^{-1}$ existe. La stabilité est donc garantie par la normalisation, sur tout
le domaine admissible, et non supposée. Numériquement, $\rho(W)=0{,}562\,g$. On divise par un
scalaire et non ligne à ligne, ce qui préserve l'hétérogénéité de ce que \emph{reçoit} chaque
pilier (P1 reçoit trois fois moins que P2) : la revendication \og P1 émet fort et reçoit peu\fg{}
survit à la normalisation.

% SOURCES-SCRIPTS: 03

\section{La cascade est un processus de branchement}
\label{sec:branch}

Le modèle simultané sur latentes n'est pas celui que les données pourraient identifier. On lui
substitue une contagion retardée sur incidents, qui est un processus de
branchement multitype : un incident sur le pilier~$k$ engendre, à la période suivante, un
nombre aléatoire d'incidents sur les autres piliers. Toute la théorie de l'extinction des
épidémies s'applique, et l'objet central n'est plus $W$ mais la matrice de génération
suivante $M$, qui compte des incidents et non des écarts-types latents :
\begin{equation}
M_{jk}\;=\;\Phi(\text{base}_j+W_{jk})-\Phi(\text{base}_j).
\end{equation}

\begin{cle}
\textbf{Deux rayons spectraux, à ne pas confondre.} $\rho(W)$ gouverne le système latent
simultané (il doit être $<1$ pour que $(I-W)^{-1}$ existe) ; $R_0=\rho(M)$ gouverne la cascade
d'incidents, qui s'éteint si et seulement si $R_0<1$ (le $R_0$ des épidémiologistes). Avec une
incidence de fond de $4{,}5\,\%$ et $g=0{,}9$, on obtient $\rho(W)=0{,}506$ et
$R_0=0{,}054$ : les deux lectures sont stables, mais pas du même ordre.
\end{cle}

\begin{figure}[htbp]\centering
\figover{K3_branchement_R0.png}
\caption{(a) Sur le domaine admissible $g\le1$, les deux rayons spectraux restent $<1$ ; les
seuils critiques ($\rho(W)=1$ à $g=1{,}78$ ; $R_0=1$ à $g=8{,}1$) sont hors domaine. (b) La
progéniture $(I-M)^{-1}$ reproduit exactement le classement ROOT du jugement d'expert.}
% SOURCES-SCRIPTS: 03
\end{figure}

Deux résultats en découlent. D'une part, la forme simultanée est structurellement
alarmiste : même normalisée, elle approche son seuil critique plus de \emph{quatre} fois plus
vite que la cascade réelle ($g^\star=1{,}78$ contre $8{,}1$). Une lecture naïve de TRANS crierait à
l'emballement systémique dans un régime où la contagion s'éteint tranquillement. D'autre part,
la progéniture $(I-M)^{-1}$ se lit comme le nombre de descendants d'un incident initial : le
calcul donne $0{,}109$ (P1), $0{,}069$ (P4), $0{,}058$ (P2), $0{,}054$ (P3), $0{,}026$ (P5), un
classement identique à celui du barème ROOT (Spearman $=1{,}00$), et ce pour toute
valeur de $g$.

\begin{attention}
\textbf{Cohérence interne, et non validation externe.} ROOT et TRANS sont issus du jugement du
même expert ; retrouver ROOT à partir de TRANS n'est pas une validation externe mais un test de
\emph{cohérence interne} du jugement, réussi. Sa portée réelle est une réduction de
paramètres : ROOT se déduit de TRANS au lieu d'être posé, soit cinq paramètres d'expert en
moins.
\end{attention}

% SOURCES-SCRIPTS: 03 29

\section{Le seuil $K_j$ : ce qui remplace $\Phi^{-1}(\mathrm{PD})$}

Il faut distinguer deux seuils vivant dans deux espaces : $u_j$, le seuil de matérialité
(DORA art.~18), dans l'espace des sévérités observables (euros, clients, heures) ; et $K_j$, le
seuil sur la variable latente, sans unité. On ne peut pas prendre l'un pour l'autre : on
apparie leurs probabilités, en choisissant $K_j$ pour que l'événement latent
$\{X_{ij}\ge K_j\}$ ait la même probabilité que l'incident majeur observable
$\{S_{ij}\ge u_j\}$ :
\begin{equation}
\label{eq:Kj}
\boxed{\;K_j\;=\;F^{-1}(1-p_j), \qquad p_j=\mathbb P(S_{ij}\ge u_j).\;}
\end{equation}

\begin{cle}
\textbf{La littérature fait déjà exactement cela.} Dans la représentation CreditMetrics des
matrices de transition \citep{jpmorgan1997}, le seuil de migration vaut
$L=\Phi^{-1}(\text{fréquence cumulée observée})$ : on inverse \emph{toujours} une répartition. Ce
qu'on reproche à $\Phi^{-1}(\mathrm{PD})$ n'est pas l'inversion, ce sont ses deux entrées : une
fréquence non mesurable (remplacée par $p_j$, définie par le règlement, donc stable et
comparable entre entités) et une queue gaussienne trop fine (remplacée par une loi $F$ à queue
lourde). La probabilité de dépassement devient alors une sortie du modèle,
conditionnelle à l'état systémique.
\end{cle}

\noindent Le biais sur $K_j$ ne vient pas de la méthode d'estimation de la queue mais de la
\textbf{sous-déclaration} : en notant $q(s)$ la probabilité qu'un incident de sévérité $s$ soit
déclaré, on a l'identité exacte $\hat p_j^{\text{naïf}}/p_j=\mathbb E[q(S)\mid S\ge u_j]<1$. On
sous-estime $p_j$, donc on sur-estime $K_j$, donc on sous-provisionne, et ce biais s'effondre
quand la barre monte dans la queue ($+0{,}369$ à $u=0{,}5$ ; $+0{,}001$ à $u=40$). La
barre DORA est défendable parce qu'elle est haute, pas parce qu'elle est réglementaire, et
c'est une propriété vérifiable : il suffit d'estimer $q(u_j)$. Le rôle de l'EVT n'est pas de
débiaiser cette fréquence (extrapoler le taux depuis un seuil haut n'est pas pilotable), mais de
modéliser la sévérité au-dessus de la barre, donc le capital ($-0{,}8\,\%$ de biais
contre $-63{,}9\,\%$ pour la lognormale).

% SOURCES-SCRIPTS: 04 02

\section{Calibration et faisabilité}
\label{sec:calib-cascade}

L'asymétrie de $W$ ne s'identifie que par l'information temporelle (qui déclenche avant
qui) : le modèle de calibration est dynamique, et l'estimateur est littéralement un
probit (le modèle générateur a un bruit gaussien et un seuil), non un logit qui gonflerait les
magnitudes d'un facteur $2{,}1$. Sur données \emph{simulées}, ce protocole fonctionne : le probit
retrouve $W$ dans ses unités (corrélation $0{,}92$, pente $1{,}08$). Le protocole n'est donc pas
en cause dans ce qui suit.

\begin{attention}
\textbf{Ce que ce protocole donne sur données réelles : rien.} Appliqué aux sources
disponibles, il échoue, et il échoue pour des raisons qui tiennent aux \emph{données} et non à
la méthode. Le chapitre~\ref{ch:identifiabilite} le démontre sur trois sources qui échouent pour
des raisons distinctes et en tire la frontière d'identification ; le
chapitre~\ref{chap:partielle} en tire les conséquences sur le capital. Une précision de méthode
s'impose dès ici : le test intuitif du \og temps inversé \fg{} est dégénéré
sur des comptages, puisque la matrice obtenue en renversant la chronologie vaut, par
construction, la transposée. Il ne prouve rien, et seul le placebo par permutation est valide.
\end{attention}

Le présent chapitre en donne la substance ; le détail du banc d'essai
$\Phi^{-1}$ contre EVT, de la calibration du facteur systémique et des diagnostics de
puissance est reporté en annexe.

% SOURCES-SCRIPTS: 03 05

\section{Formalisation : définitions et propriétés}
\label{sec:formalisation}

Cette section consolide le modèle sous une forme opposable à un lecteur exigeant. On note
$P=\{1,\dots,5\}$ les piliers, $W\in[0,1]^{P\times P}$ la matrice de contagion avec
$W_{jj}=0$, et $s_{\max}=\max_j\sum_k \mathrm{TRANS}_{jk}$. Les démonstrations sont
rassemblées à l'annexe~\ref{ann:cascade}.

\begin{definition}[Cascade indépendante]
\label{def:cascade}
À chaque couple ordonné $(j,k)$, $j\neq k$, on associe une variable de Bernoulli
indépendante $B_{jk}\sim\mathcal B(W_{jk})$. Le \emph{graphe d'arêtes vivantes} $G$ possède
l'arc $j\to k$ si et seulement si $B_{jk}=1$. Pour un pilier d'amorce $a$, l'\emph{ensemble
atteint} est
\[
  S_a \;=\; \{a\}\cup\{k\in P : \text{il existe un chemin dirigé } a\rightsquigarrow k \text{ dans } G\}.
\]
C'est la construction standard de la cascade indépendante \citep{kempe2003}, équivalente à une percolation
de liens sur le graphe dirigé pondéré par $W$.
\end{definition}

\begin{proposition}[Forme réduite et descendance espérée]
\label{prop:reduite}
Le nombre espéré de marches dirigées de longueur $n$ de $a$ vers $k$ vaut $(W^n)_{ak}$, et la
descendance totale espérée, comptée avec multiplicité, est
\[
  \sum_{n\ge 0}(W^n)_{ak} \;=\; \bigl((I-W)^{-1}\bigr)_{ak},
\]
série convergente si et seulement si $\rho(W)<1$. L'ensemble atteint étant compté sans
multiplicité, on a la majoration $\mathbb E\,|S_a|\le \bigl[(I-W)^{-1}\mathbf 1\bigr]_a$, et
l'écart mesure exactement les piliers atteints par plusieurs chemins.
\end{proposition}

\begin{proposition}[Sous-criticité par normalisation de Leontief]
\label{prop:leontief}
Sous la normalisation~\eqref{eq:norm}, $W=(g/s_{\max})\,\mathrm{TRANS}$, la matrice $W$ est
positive et chacune de ses lignes somme à au plus $g$ ; par conséquent
\[
  \rho(W)\;\le\;\max_j\sum_k W_{jk}\;\le\;g\;<\;1 .
\]
La convergence de $(I-W)^{-1}$, donc la stabilité de la cascade, est ainsi \emph{garantie sur
tout le domaine admissible}, et non supposée.
\end{proposition}

\begin{proposition}[Loi exacte de l'ensemble atteint]
\label{prop:loi}
Pour tout $A\subseteq P$ avec $a\in A$,
\[
  \mathbb P(S_a=A)\;=\;R(A)\,B(A),\qquad
  B(A)=\!\!\prod_{j\in A,\;k\notin A}\!\!(1-W_{jk}),
\]
où $B(A)$ est la probabilité qu'aucune arête ne fuie hors de $A$, et $R(A)$, la probabilité
que $a$ atteigne \emph{tout} $A$ par les seules arêtes internes, vérifie $R(\{a\})=1$ et
\[
  R(A)\;=\;1-\sum_{\substack{A'\subsetneq A\\ a\in A'}} R(A')\!\!\prod_{j\in A',\;k\in A\setminus A'}\!\!(1-W_{jk}).
\]
\end{proposition}

\noindent Le moteur qui produit les capitaux publiés n'emploie pas cette cascade mais sa variante
à un successeur par pas, la marche auto-évitante : depuis le pilier courant $j$, la cascade passe
au pilier $k$ avec la probabilité $W_{jk}$, et s'arrête sinon, ou dès qu'elle revient sur un
pilier déjà touché. Les deux constructions ont la même matrice moyenne $W$ et le même nombre
moyen de descendants directs, et la loi de l'ensemble atteint se calcule exactement dans les deux
cas, sans simulation. Elles diffèrent par la forme de la loi de taille et non par le capital :
passer de la marche au branchement déplace le SCR de $-4\,\%$. Les
propositions~\ref{prop:reduite} à~\ref{prop:mono} sont énoncées pour la cascade indépendante,
qui en donne la forme la plus simple ; pour la marche, la croissance du capital en $g$ est
vérifiée numériquement au \S\ref{sec:invariance-g}.

\begin{proposition}[Monotonie stochastique]
\label{prop:mono}
Si $W\le W'$ terme à terme, alors $S_a(W)$ est stochastiquement dominé par $S_a(W')$. Toute
fonctionnelle croissante de $S_a$, en particulier la perte agrégée et sa VaR, est donc
croissante en chaque entrée de $W$. En revanche une perturbation \emph{antisymétrique}
$A$ déplace $W_{jk}$ et $W_{kj}$ en sens opposés : le capital n'est pas monotone en $A_{jk}$,
ce qui justifie de balayer les sommets de l'ensemble admissible
(chapitre~\ref{chap:partielle}) sans revendiquer de bornes exactes.
\end{proposition}

\begin{proposition}[Dépendance à l'ordre]
\label{prop:ordre}
La vraisemblance d'une chaîne ordonnée $\ell(i_1,\dots,i_n)=\mathrm{ROOT}_{i_1}\prod_k
\mathrm{TRANS}_{i_k i_{k+1}}$ est invariante par permutation si et seulement si
$\mathrm{TRANS}$ est symétrique. Comme $\mathrm{TRANS}$ est asymétrique, la criticité
\emph{n'est pas une fonction de l'ensemble} des piliers touchés : aucune copule ni aucun
score additif, invariants par permutation, ne peut la reproduire.
\end{proposition}

\begin{proposition}[Ce que la co-occurrence voit de la direction]
\label{prop:nonid}
Soit la décomposition unique $W=S+A$, avec $S=\tfrac12(W+W^\top)$ et $A=\tfrac12(W-W^\top)$.
\emph{(i)} Sous une amorce uniforme, de probabilité $r$ pour chaque pilier, la probabilité qu'un
même sinistre touche à la fois $j$ et $k$ vaut
\[
  \sum_{a\in P} r\,\mathbb P\bigl(\{j,k\}\subseteq S_a\bigr)
  \;=\; r\,(W_{jk}+W_{kj})+O\bigl(\lVert W\rVert^2\bigr)
  \;=\; 2r\,S_{jk}+O\bigl(\lVert W\rVert^2\bigr) :
\]
au premier ordre, la co-occurrence ne dépend de $W$ que par $S$, et $W$ et $W^\top$ y sont
indiscernables. \emph{(ii)} Au-delà du premier ordre, ou sous une amorce non uniforme, la loi de
l'ensemble atteint dépend de $A$ en général : $W$ et $W^\top$ ne sont pas observationnellement
équivalents pour la loi complète des ensembles de piliers touchés. \emph{(iii)}
$\mathrm{SCR}(W)\neq\mathrm{SCR}(W^\top)$ en général.
\end{proposition}

% SOURCES-SCRIPTS: 32 61 29 98

\section{Vérification numérique des trois propriétés qui portent l'argument}
\label{sec:verif-proprietes}

Parmi les six énoncés ci-dessus, trois portent l'argument du mémoire. Chacun est accompagné de
sa vérification numérique : l'énoncé
% SOURCES-SCRIPTS: 32
dit ce qui doit se produire, le calcul montre que cela se produit.

\paragraph{La criticité dépend de l'ordre (proposition~\ref{prop:ordre}).} Sur les $26$
sous-ensembles de taille au moins deux, $22$ changent de criticité selon l'ordre, avec
un écart atteignant $6$ points sur une échelle de $10$, soit $60\,\%$ de l'échelle, pour
un \emph{même} ensemble de piliers. Le contre-exemple minimal suffit à le voir : sur la paire
$\{P_1,P_2\}$, le classeur donne $\mathrm{TRANS}[1][2]=0{,}80$ et $\mathrm{TRANS}[2][1]=0{,}20$,
d'où des criticités de $8$ et $2$ pour le même ensemble.

\begin{cle}
\textbf{La contraposée est l'argument utile.} Toute mesure de risque qui ne dépend que de
l'ensemble touché est invariante par permutation. La proposition~\ref{prop:ordre} exhibe donc un
objet qu'\emph{aucune} copule et qu'\emph{aucun} score additif ne reproduit, et l'hypothèse
requise est minimale : il suffit que $\mathrm{TRANS}$ soit asymétrique, sans qu'aucune
\emph{valeur} de la direction ne soit revendiquée.
\end{cle}

\begin{cle}
\textbf{Ce qui n'est pas revendiqué, et pourquoi le dire.} Une propriété plus forte, la
\emph{non-transitivité} de la dominance entre piliers, a été testée dans trois formulations :
la transitivité des corrélations qu'impose un modèle à facteur unique, la positivité de la
dépendance symétrisée, et l'existence de cycles de dominance par paire. Aucune ne se
vérifie : zéro violation de l'inégalité de transitivité sur $420$ triplets, valeur propre
minimale $+0{,}405$ donc matrice de corrélation parfaitement valide, et zéro cycle. Cette
revendication est donc retirée. La dépendance à l'ordre exclut exactement les mêmes
concurrents et, elle, se démontre.
\end{cle}

\paragraph{La borne de propagation tient, et son écart se lit (propositions~\ref{prop:reduite}
et~\ref{prop:leontief}).} Avec le classeur actuel, $\rho(W)=0{,}506$. La majoration
$\E\,|S_a|\le[(I-W)^{-1}\mathbf 1]_a$ est vérifiée sur les cinq amorces et sur $300$ matrices de
l'ensemble admissible, sans aucun contre-exemple. L'écart entre la borne et la valeur exacte va
de $10{,}6\,\%$ (amorce P5) à $19{,}9\,\%$ (amorce P1). Cet écart n'est pas une
imprécision : il mesure exactement les \emph{collisions}, c'est-à-dire les piliers atteints par
plusieurs chemins, que la forme linéaire compte plusieurs fois. Sur cinq piliers fortement
connectés il n'est pas négligeable, ce qui justifie de conserver le calcul exact par
sous-ensembles (proposition~\ref{prop:loi}) plutôt que la seule forme close.

\paragraph{La direction échappe aux statistiques disponibles, et cela coûte
(proposition~\ref{prop:nonid}).} Le test directionnel ne rejette pas l'hypothèse $A=0$
(chapitre~\ref{ch:identifiabilite} : asymétrie observée $119$ contre un nul de permutation à
$128 \pm 27$, soit $z=-0{,}33$, et aucune arête significative sur $18$). Retourner la matrice
déplace pourtant le capital de $3{,}6\,\%$ et la perte moyenne de $7{,}9\,\%$ (les niveaux sont
établis en partie~\ref{part:resultats} ; seuls les écarts importent ici), et surtout la priorité
de remédiation bascule de P1 à P4. Sur l'ensemble admissible, l'écart entre une matrice et sa
transposée atteint $1\,173$~M\euro{}.

La loi complète des ensembles touchés, elle, distingue les deux matrices, ce que la partie (ii) de
la proposition énonce. Sur la matrice publiée, avec l'amorce et la marche du moteur, un sinistre
touche en moyenne $1{,}931$ pilier sous $W$ et $1{,}746$ sous $W^\top$, la part des sinistres
multi-piliers passe de $62{,}31\,\%$ à $50{,}77\,\%$, et la probabilité que P2 et P5 cèdent
ensemble tombe de $0{,}127$ à $0{,}066$. Sous amorce uniforme, l'écart entre les deux matrices
s'efface avec l'intensité de la contagion : il ne vaut plus que $0{,}86\,\%$ quand $W$ est ramenée
au cinquantième de sa valeur, ce qui est la partie (i).

\begin{cle}
\textbf{L'énoncé à retenir : indiscernables pour les statistiques disponibles, matériellement
différentes.} C'est la formulation exacte de la frontière d'identification, et elle ne repose sur
\emph{aucun} dire d'expert. Elle chiffre en euros de capital, et en inversion de priorité, le
prix de ne pas savoir dans quel sens la contagion circule. Elle désigne aussi la donnée qui
réduirait cette ignorance : la direction n'est pas inaccessible par principe, elle est absente
des sources, et un registre qui consignerait l'ensemble des piliers touchés par chaque incident
en porterait une partie. C'est ce qui fonde à la fois le modèle en couches
(chapitre~\ref{chap:partielle}) et la spécification de reporting adressée au régulateur.
\end{cle}

% =====================================================================

% SOURCES-SCRIPTS: 40 03 98
